%%%%%%%%%%%%%%%%%%%%%%%%%%%%%%%%%%%%%%%%%
% Medium Length Professional CV
% LaTeX Template
% Version 2.0 (8/5/13)
%
% This template has been downloaded from:
% http://www.LaTeXTemplates.com
%
% Original author:
% Trey Hunner (http://www.treyhunner.com/)
%
% Important note:
% This template requires the resume.cls file to be in the same directory as the
% .tex file. The resume.cls file provides the resume style used for structuring the
% document.
%
%%%%%%%%%%%%%%%%%%%%%%%%%%%%%%%%%%%%%%%%%

%----------------------------------------------------------------------------------------
%	PACKAGES AND OTHER DOCUMENT CONFIGURATIONS
%----------------------------------------------------------------------------------------


\documentclass[
	%a4paper, % Uncomment for A4 paper size (default is US letter)
	% 11pt, % Default font size, can use 10pt, 11pt or 12pt
]{resume} % Use the resume class

% \usepackage{ebgaramond} % Use the EB Garamond font
\usepackage{newpxtext}
\usepackage{hyperref}
\usepackage{enumitem}
\hypersetup{colorlinks=true}


\name{\underline{Sun} Guangda} % Your name
% \address{National University of Singapore \\ 21 Lower Kent Ridge Rd, Singapore 119077} % Your address

\address{\href{https://github.com/sgdxbc}{GitHub: sgdxbc} \\ \href{mailto:guangda@u.nus.edu}{guangda@u.nus.edu}} % Your phone number and email

\begin{document}

Ph.D. in Computer Science from NUS, now an R\&D engineer at ByteDance working on large-scale data infrastructure.
I design, build, and evaluate distributed systems end to end, mostly in Rust:
Byzantine fault-tolerant replication, blockchain and decentralized storage, and high-performance network systems.
I also have a background in compilers and constraint solving.
Published at SIGCOMM, NSDI, CoNEXT, and IEEE/ACM ToN.
I enjoy turning research ideas into reliable, high-performance systems.

%----------------------------------------------------------------------------------------
%	EDUCATION SECTION
%----------------------------------------------------------------------------------------
\begin{rSection}{Education}
    \textbf{National University of Singapore} \hfill {\em 2020 - 2026} \\
    Ph.D. in Computer Science \\
    Advisor: Jialin Li, Assistant Professor \\

    \textbf{Xi'an Jiaotong University}  \hfill {\em 2014 - 2020} \\
    % School of Electronic and Information Engineering \\
    Bachelor of Engineering in Computer Science and Technology, \emph{Honors Youth Program} \\
    Advisor: Hao Li, Associate Professor
\end{rSection}

%----------------------------------------------------------------------------------------
%	INDUSTRY EXPERIENCE SECTION
%----------------------------------------------------------------------------------------
\begin{rSection}{Industry Experience}
    \begin{rSubsection}{ByteDance Pte. Ltd.}{Aug 2026 - present}{Research \& Development}{Singapore}
        \item Develop and optimize the data path of the streaming learning-sample pipeline, serving recommendation (e.g., Douyin \& TikTok video feeds), advertising, and search.
        \item Optimize storage in large-scale clusters (10k+ cores); improved network communication to retire a message queue that consumed several TBs of storage.
    \end{rSubsection}
\end{rSection}

%----------------------------------------------------------------------------------------
%	RESEARCH EXPERIENCE SECTION
%----------------------------------------------------------------------------------------
\begin{rSection}{Research Experience}
    \begin{rSubsection}{BigBFT: scaling BFT via state sharding}{May 2025 - present}{APSys '26, first author}{}
        \item Designed a BFT transactional store that shards state but not execution: one consensus protocol orders all transactions while node shards serve state partitions in parallel, without weakening fault tolerance.
        \item Built the storage layer with tiered updates, state snapshots, authenticated data structures, and erasure coding; implemented in Rust and evaluated with geo-distributed deployments on AWS.
    \end{rSubsection}

    \begin{rSubsection}{Decentralized storage network with scalable volume and security}{Jun 2022 - present}{Under submission}{}
        \item Designed a decentralized storage network that achieves data availability, storage scalability, and blockchain efficiency simultaneously, using verifiable sampling-based placement and verifiable rateless erasure coding.
            Implemented in Rust; evaluated with discrete-event simulation and geo-distributed deployments on AWS.
    \end{rSubsection}

    \begin{rSubsection}{NeoBFT: BFT replication with authenticated in-network ordering}{Jan 2021 - Nov 2023}{SIGCOMM '23, first author}{}
        \item Proposed an authenticated in-network ordering primitive and built a high-performance BFT protocol on it.
        \item Implemented the primitive in P4 on Tofino switches and the protocol in Rust; evaluated at scale on AWS.
    \end{rSubsection}

    \begin{rSubsection}{LemonNFV: secure, high-performance network function virtualization}{Jun 2020 - May 2022}{NSDI '23}{}
        \item Co-implemented a framework that consolidates heterogeneous network functions at line speed, isolated with Intel MPK and a custom dynamic loader on DPDK; wrote stackful coroutines in C and x86 assembly.
    \end{rSubsection}

    \begin{rSubsection}{Rubik: programming network stacks for middleboxes}{Aug 2018 - Apr 2020}{NSDI '21, IEEE/ACM ToN}{}
        \item Implemented a Python compiler that generates optimized C network stacks from short protocol specifications.
    \end{rSubsection}

    \begin{rSubsection}{Solver-based compiler for network programming languages}{Feb 2018 - Feb 2019}{CoNEXT '20, IEEE/ACM ToN}{}
        \item Implemented a compiler in Cython that compiles programs written in multiple software-defined networking languages into mixed-integer programming problems, solves them with Gurobi, and translates the solutions into global network configurations.
        \item Designed parallel problem simplification and pre-splitting, cutting solving time from weeks to minutes.
    \end{rSubsection}
\end{rSection}

%----------------------------------------------------------------------------------------
%	PUBLICATIONS SECTION
%----------------------------------------------------------------------------------------
\begin{rSection}{Publications}
\begin{enumerate}[label={[}\arabic*{]}]
% \setlength\itemsep{-0.5em}

\item
    \textbf{Guangda Sun}, Jialin Li.
    \textit{BigBFT: Scaling BFT without Compromising Fault Tolerance via State Sharding.}
    ACM SIGOPS Asia-Pacific Workshop on Systems (APSys '26).
    \href{https://doi.org/10.1145/3838177.3841729}{paper},
    \href{https://github.com/sgdxbc/sgdxbc/blob/main/papers/bigbft_apsys26.pdf}{pdf}

\item
    \textbf{Guangda Sun}, Mingliang Jiang, Xin Zhe Khooi, Yunfan Li, Jialin Li.
    \textit{NeoBFT: Accelerating Byzantine Fault Tolerance Using Authenticated In-Network Ordering.}
    Proceedings of the ACM SIGCOMM 2023 Conference.
    \href{https://doi.org/10.1145/3603269.3604874}{paper},
    \href{https://github.com/sgdxbc/sgdxbc/blob/main/papers/3603269.3604874.pdf}{pdf},
    \href{https://github.com/nus-sys/neobft-artifact}{code}


\item
    Hao Li, Yihan Dang, \textbf{Guangda Sun}, Changhao Wu, Peng Zhang, Danfeng Shan, Tian Pan, and Chengchen Hu.
    \textit{Programming Network Stack for Physical Middleboxes and Virtualized Network Functions.}
    IEEE/ACM Transactions on Networking 32.2 (2023).
    \href{https://ieeexplore.ieee.org/document/10273850}{paper},
    \href{https://github.com/sgdxbc/sgdxbc/blob/main/papers/TNET.2023.3307641.pdf}{pdf}

\item
    Hao Li, Yihan Dang, \textbf{Guangda Sun}, Guyue Liu, Danfeng Shan, Peng Zhang.
    \textit{LemonNFV: Consolidating Heterogeneous Network Functions at Line Speed.}
    20th USENIX Symposium on Networked Systems Design and Implementation (NSDI '23).
    \href{https://www.usenix.org/conference/nsdi23/presentation/li-hao}{paper},
    \href{https://github.com/sgdxbc/sgdxbc/blob/main/papers/nsdi23-li-hao.pdf}{pdf},
    \href{https://github.com/ants-xjtu/neptune}{code}

\item
    Hao Li, Changhao Wu, \textbf{Guangda Sun}, Peng Zhang, Danfeng Shan, Tian Pan, Chengchen Hu.
    \textit{Programming Network Stack for Middleboxes with Rubik.}
    18th USENIX Symposium on Networked Systems Design and Implementation (NSDI '21).
    \href{https://www.usenix.org/conference/nsdi21/presentation/li}{paper},
    \href{https://github.com/sgdxbc/sgdxbc/blob/main/papers/nsdi21-li.pdf}{pdf},
    \href{https://github.com/ants-xjtu/rubik}{code}

\item
    Hao Li, Peng Zhang, \textbf{Guangda Sun}, Wanyue Cao, Chengchen Hu, Danfeng Shan, Tian Pan, Qiang Fu.
    \textit{Compiling Cross-Language Network Programs Into Hybrid Data Plane.}
    IEEE/ACM Transactions on Networking 30.3 (2021).
    \href{http://ieeexplore.ieee.org/document/9648264}{paper},
    \href{https://github.com/sgdxbc/sgdxbc/blob/main/papers/coder-ton22-li.pdf}{pdf}

\item
    Hao Li, Peng Zhang, \textbf{Guangda Sun}, Chengchen Hu, Danfeng Shan, Tian Pan, Qiang Fu.
    \textit{A Modular Compiler for Network Programming Languages.}
    Proceedings of the 16th International Conference on emerging Networking EXperiments and Technologies (CoNEXT '20).
    \href{https://dl.acm.org/doi/10.1145/3386367.3432063}{paper},
    \href{https://github.com/sgdxbc/sgdxbc/blob/main/papers/coder-conext20-li.pdf}{pdf},
    \href{https://github.com/ants-xjtu/sidfam-v3}{code}

\item
    Hao Li, Peng Zhang, \textbf{Guangda Sun}, Chengchen Hu, Danfeng Shan, Tian Pan, Qiang Fu.
    \textit{An Intermediate Representation for Network Programming Languages.}
    Proceedings of the 4th Asia-Pacific Workshop on Networking (APNet '20).
    \href{https://dl.acm.org/doi/10.1145/3411029.3411030}{paper},
    \href{https://github.com/sgdxbc/sgdxbc/blob/main/papers/apnet20-final4.pdf}{pdf}

\end{enumerate}
\end{rSection}

%----------------------------------------------------------------------------------------
%	SELECTED PROJECTS SECTION
%----------------------------------------------------------------------------------------
\begin{rSection}{Selected Projects}
\begin{tabular}{ @{} >{\bfseries}l p{0.8\textwidth} }
\href{https://github.com/neatsys/neatworks}{Neatworks} & Reusable Rust building blocks and reference implementations for distributed protocols \\
\href{https://github.com/sgdxbc/izuko}{izuko} & A profiler and monitor for IPFS \\
\href{https://github.com/whoiscc/simpile}{simpile} & A memory allocator written in Rust and fuzz-tested \\
\href{https://github.com/whoiscc/shattuck}{Shattuck} & A scripting language interpreter written in Rust
\end{tabular}
\end{rSection}

%----------------------------------------------------------------------------------------
%	TECHNICAL STRENGTHS SECTION
%----------------------------------------------------------------------------------------
\begin{rSection}{Technical Strengths}
\begin{tabular}{ @{} >{\bfseries}l p{0.71\textwidth} }
Programming languages
    & Daily using: Rust, Python, C++, Java, \LaTeX \\
    & Have completed projects with: C, TypeScript, P4, x86 assembly, Coq \\
    & Familiar with: Go, OCaml, Haskell, Erlang, C\#, Ruby \\
Areas
    & Distributed systems, BFT, storage, high-performance networking, solvers \\
Tools \& platforms
    & Linux, AWS, Terraform, Cloudflare, DPDK, Tofino, Gurobi, Claude Code, Codex
\end{tabular}
\end{rSection}

%----------------------------------------------------------------------------------------
%	HONORS & TEACHING SECTION
%----------------------------------------------------------------------------------------
\begin{rSection}{Honors \& Teaching}
\begin{tabular}{ @{} >{\bfseries}ll }
2023 & Research Achievement Award, School of Computing, NUS \\
2020 & NUS Research Scholarship \\
Teaching & CS5223 Distributed Systems, NUS \\
\end{tabular}
\end{rSection}

\iffalse
\begin{rSection}{Projects For Fun}
    VS Code extensions:
    \begin{itemize}
        \item \href{https://github.com/whoiscc/rust-comment}{rust-comment}: A VS Code extension for correct comment width adhere to \href{https://doc.rust-lang.org/style-guide/\#comments}{The Rust Style Guide}.
        \item \href{https://github.com/whoiscc/shift-saver}{Shift Saver}: A VS Code extension reducing the use of shift key when naming.
    \end{itemize}

    Languages:
    \begin{itemize}
        \item \href{https://github.com/whoiscc/shattuck}{Shattack}: A scripting language interpreter written in Rust.
        \item \href{https://github.com/whoiscc/cgen}{cgen}: A code generator in Python as a generic programming solution for C.
    \end{itemize}

    Memory allocators:
    \begin{itemize}
        \item \href{https://github.com/sgdxbc/MinimalHeap}{MinimalHeap}: A well-tested baseline implementation.
        \item \href{https://github.com/ants-xjtu/memgrid}{MemGrid}: A memory allocator that supports allocating in pre-allocated buffer.
        \item \href{https://github.com/whoiscc/simpile}{simpile}: A memory allocator written in Rust and fuzz-tested.
    \end{itemize}

    \href{https://github.com/neatsys/neatworks}{Neatworks}: A researching codebase for distributed protocols in Rust.
    Reusable network transportation module and reference implementations for various seminar works.

    \href{https://www.zhihu.com/people/cowsay/posts}{More than 100 articles on Zhihu}. 4.4k up votes in total.

    \href{https://littlebot.blog}{Personal blog} (open source soon).
    A content-focused minimal website generated directly from Python scripts without static site generator.
    Deployed with Cloudflare Pages and RAG-based searching backed by Cloudflare AI Search.

    \href{https://github.com/whoiscc/fastray}{fastray}: My Rust implementation of \href{https://raytracing.github.io/}{Ray Tracing in One Weekend}.

    \href{https://github.com/whoiscc/twenty48}{twenty48}: A 2048 game with perfect animation for arbitrarily frequent operations.

    \href{https://github.com/sgdxbc/izuko}{izuko}: A profiler and monitor for IPFS.

    Progress management application (public available soon).
    Frontend with Vue and backend with Hono and SQLite.
    Deployed with Cloudflare Workers, D1 and R2.
\end{rSection}
\fi

%----------------------------------------------------------------------------------------

\end{document}
